% Luc Destombe --- licence CC BY-NC-SA 4.0
% https://creativecommons.org/licenses/by-nc-sa/4.0/deed.fr
% Fichier autonome : compiler avec pdflatex, deux fois (signets, nombre de pages).
\documentclass[11pt,a4paper]{article}
\makeatletter
\newif\ifeleve
\elevefalse

\newif\iflycee@palatino
\lycee@palatinotrue

\RequirePackage[utf8]{inputenc}
\RequirePackage[T1]{fontenc}
\RequirePackage[french]{babel}
\RequirePackage{etoolbox}

\newcommand{\ShorthandsOff}{\@ifundefined{shorthandoff}{}{\shorthandoff{;:!?}}}
\newcommand{\ShorthandsOn}{\@ifundefined{shorthandon}{}{\shorthandon{;:!?}}}
\ShorthandsOff
\AtBeginDocument{\ShorthandsOn}
\AtBeginEnvironment{tikzpicture}{\ShorthandsOff}

\RequirePackage{geometry}
\RequirePackage{fancyhdr}
\RequirePackage{lastpage}
\RequirePackage{multicol}
\RequirePackage{array}
\RequirePackage{enumitem}
\RequirePackage{xcolor}
\RequirePackage{needspace}

\RequirePackage{amsmath,amssymb}
\RequirePackage{mathpazo}
\DeclareMathSizes{10.95}{10.95}{8}{5.5}
\begingroup\catcode`\_=\active \catcode`\^=\active
\gdef_#1{\sb{\scriptscriptstyle #1}}
\gdef^#1{\sp{\scriptscriptstyle\lycee@moinsepais #1}}
\endgroup
\newcommand\lycee@moins{\mathbin{%
  \setbox\z@\hbox{$\m@th\scriptscriptstyle\lycee@moinsorig$}%
  \dimen@=.55\wd\z@ \advance\dimen@-.5pt
  \hbox to.75\wd\z@{\hss$\m@th\scriptscriptstyle\vcenter{\hbox{%
    \tikz\draw[line width=.5pt,line cap=round](0,0)--(\the\dimen@,0);}}$\hss}}}
\begingroup\lccode`\~=`\- \lowercase{\endgroup\let~}\lycee@moins
\newcommand\lycee@moinsepais{\mathcode`\-="8000 }
\AtBeginDocument{\mathchardef\lycee@moinsorig=\mathcode`\-\relax}
\AtBeginDocument{\catcode`\_=12 \mathcode`\_="8000
  \catcode`\^=12 \mathcode`\^="8000 }
\newcommand\lycee@exposants{%
  \fontdimen13\textfont2=.48\fontdimen6\textfont2
  \fontdimen14\textfont2=.48\fontdimen6\textfont2
  \fontdimen15\textfont2=.48\fontdimen6\textfont2
  \fontdimen13\scriptfont2=.48\fontdimen6\scriptfont2
  \fontdimen14\scriptfont2=.48\fontdimen6\scriptfont2
  \fontdimen15\scriptfont2=.48\fontdimen6\scriptfont2}
\every@math@size\expandafter{\the\every@math@size\lycee@exposants}
\newlength{\retraitcalculs}
\setlength{\retraitcalculs}{2em}

\RequirePackage{tikz}
\usetikzlibrary{arrows.meta,calc,babel,shapes.misc}
\RequirePackage[most]{tcolorbox}
\tcbuselibrary{skins,breakable}

\definecolor{coulDef}{HTML}{1F4E79}
\definecolor{coulProp}{HTML}{7030A0}
\definecolor{coulRem}{HTML}{C55A11}
\definecolor{coulEx}{HTML}{2E7D32}
\definecolor{coulExo}{HTML}{B03A2E}
\definecolor{coulPart}{HTML}{1F4E79}
\definecolor{coulMeth}{HTML}{1B6E4F}
\definecolor{coulCourbe}{HTML}{0B5ED7}
\definecolor{coulCourbeB}{HTML}{C00000}
\definecolor{coulCourbeC}{HTML}{2E7D32}
\definecolor{coulGrille}{HTML}{8C8C8C}
\definecolor{coulRep}{HTML}{1B6E4F}
\definecolor{coulBar}{HTML}{2E7D32}

\newcommand{\R}{\mathbb{R}}
\newcommand{\Rs}{\mathbb{R}^{*}}
\renewcommand{\le}{\leqslant}
\renewcommand{\ge}{\geqslant}
\renewcommand{\approx}{\simeq}
\newcommand{\vect}[1]{\overrightarrow{#1}}
\newcommand{\norme}[1]{\left\lVert #1 \right\rVert}
\newcommand{\intervalle}[2]{\left[#1\,;#2\right]}
\RequirePackage{cancel}
\renewcommand{\CancelColor}{\color{coulExo}}
\newcommand{\fois}[1]{\times{\color{coulExo}#1}}
\newcommand{\barre}[1]{\cancel{#1}}

\tikzset{grille/.style={coulGrille,line width=0.4pt}}
\newcommand{\quadrillage}[4]{\draw[grille,step=1] (#1,#2) grid (#3,#4);}
\tikzset{
  croix/.style={cross out,draw,line width=0.6pt,inner sep=0pt,minimum size=4pt},
  croixdroite/.style={croix,rotate=45,line width=0.8pt},
}
\newcommand{\pt}[3]{\path (#1) node[croix]{} node[#2,font=\small,inner sep=2.5pt]{$#3$};}

\newcommand{\lycee@repere}[4]{%
  \draw[grille,step=1] (#1,#2) grid (#3,#4);
  \draw[-{Stealth[length=2mm]},thick] (#1,0)--({#3+0.4},0) node[below left=-1pt]{$x$};
  \draw[-{Stealth[length=2mm]},thick] (0,#2)--(0,{#4+0.4}) node[below left=-1pt]{$y$};
  \node[below left,font=\scriptsize,inner sep=1.5pt] at (0,0) {$0$};}
\newcommand{\lycee@gradx}[1]{\foreach \k/\l in {#1}{%
  \draw (\k,0.07)--(\k,-0.07) node[below,font=\scriptsize,inner sep=1.5pt]{$\l$};}}
\newcommand{\lycee@grady}[1]{\foreach \k/\l in {#1}{%
  \draw (0.07,\k)--(-0.07,\k) node[left,font=\scriptsize,inner sep=1.5pt]{$\l$};}}
\newcommand{\lycee@intff}[2]{\left[#1\,;#2\right]}
\newcommand{\lycee@intfo}[2]{\left[#1\,;#2\right[}
\newcommand{\lycee@intof}[2]{\left]#1\,;#2\right]}
\newcommand{\lycee@intoo}[2]{\left]#1\,;#2\right[}
\newcommand{\lycee@ens}[1]{\left\{#1\right\}}
\AtBeginDocument{%
  \providecommand{\repere}{\lycee@repere}%
  \providecommand{\gradx}{\lycee@gradx}%
  \providecommand{\grady}{\lycee@grady}%
  \providecommand{\Cf}{\mathcal{C}\sb{\scriptscriptstyle f}}%
  \providecommand{\Cg}{\mathcal{C}\sb{\scriptscriptstyle g}}%
  \providecommand{\intff}{\lycee@intff}%
  \providecommand{\intfo}{\lycee@intfo}%
  \providecommand{\intof}{\lycee@intof}%
  \providecommand{\intoo}{\lycee@intoo}%
  \providecommand{\ens}{\lycee@ens}}

\newlist{questions}{enumerate}{3}
\setlist[questions,1]{label=\arabic*\textdegree),leftmargin=*,itemsep=2pt,topsep=3pt}
\setlist[questions,2]{label=\alph*),leftmargin=*,itemsep=1pt,topsep=2pt}
\setlist[questions,3]{label=\roman*),leftmargin=*,itemsep=1pt,topsep=1pt}
\setlist[itemize]{label=\textbullet,leftmargin=*,itemsep=2pt,topsep=3pt}

\newcommand{\besoinplace}[1]{\par\penalty\@M\begingroup
  \dimen@=#1\relax\dimen@ii=\pagegoal\advance\dimen@ii-\pagetotal
  \ifdim\dimen@>\dimen@ii\ifdim\dimen@ii>\z@\vfil\fi\break\fi\endgroup}

\newcommand{\lycee@pied}{\thepage/\pageref{LastPage}}
\newcommand{\entete}[2]{%
  \pagestyle{fancy}\fancyhf{}%
  \fancyhead[L]{\footnotesize\color{coulPart}\textbf{#1}}%
  \fancyhead[R]{\footnotesize\color{coulPart}\textbf{#2}}%
  \fancyfoot[C]{\footnotesize\lycee@pied}%
  \renewcommand{\headrulewidth}{0.4pt}%
  \renewcommand{\headrule}{\hbox to\headwidth{%
    \color{coulPart!40}\leaders\hrule height \headrulewidth\hfill}}}

\RequirePackage{ccicons}
\newcommand{\lycee@licence}{{\scriptsize\color{gray}%
  \copyright\ Luc Destombe --- \ccbyncsa\ CC BY-NC-SA 4.0}}
\newif\iflycee@licence \lycee@licencetrue
\newcommand{\sanslicence}{\lycee@licencefalse}
\AtBeginDocument{\iflycee@licence\AddToHookNext{shipout/foreground}{%
  \put(0,0){\rlap{\kern\dimexpr1in+\hoffset+\oddsidemargin+\textwidth\relax
    \llap{\raisebox{-\dimexpr1in+\voffset+\topmargin+\headheight+\headsep
      +\textheight+\footskip\relax}{\lycee@licence}}}}}\fi}

\newcommand{\formulesserrees}{%
  \setlength{\abovedisplayskip}{3pt}\setlength{\belowdisplayskip}{3pt}%
  \setlength{\abovedisplayshortskip}{2pt}\setlength{\belowdisplayshortskip}{3pt}}

\setlength{\parindent}{0pt}

\RequirePackage[implicit=false,hidelinks,pdfencoding=auto,psdextra,bookmarksopen,
  bookmarksopenlevel=1,pdfcreator={},pdfproducer={}]{hyperref}
\RequirePackage{bookmark}
\hypersetup{pdfauthor={Luc Destombe},pdfkeywords={CC BY-NC-SA 4.0}}
\newcounter{lycee@signet}
\newcommand{\signet}[3][]{\stepcounter{lycee@signet}%
  \edef\lycee@dest{\if\relax\detokenize{#1}\relax
    signet.\arabic{lycee@signet}\else#1\fi}%
  \hypertarget{\lycee@dest}{}\bookmark[level=#2,dest=\lycee@dest]{#3}}
\pdfstringdefDisableCommands{%
  \def\R{R}\def\vect#1{#1}\def\Cf{Cf}\def\Cg{Cg}\def\fois#1{×#1}%
  \def\barre#1{#1}\def\intervalle#1#2{[#1;#2]}\def\intff#1#2{[#1;#2]}%
  \def\textsuperscript#1{#1}\def\dfrac#1#2{#1/#2}\def\frac#1#2{#1/#2}%
  \def\vec#1{#1⃗}}%
\geometry{top=2cm,bottom=1cm,left=1cm,right=1cm,headheight=14pt,footskip=0.6cm}

\RequirePackage{pgfplots}
\pgfplotsset{compat=1.18}
\RequirePackage{tkz-tab}

\newcounter{partiecours}
\renewcommand{\thepartiecours}{\Roman{partiecours}}
\newcommand{\partiecours}[1]{%
  \refstepcounter{partiecours}%
  \Needspace*{8\baselineskip}%
  \bigskip
  \noindent\begin{tcolorbox}[enhanced,colback=coulPart,colframe=coulPart,
      boxrule=0pt,arc=2pt,left=8pt,right=8pt,top=4pt,bottom=4pt,
      before skip=6pt,after skip=10pt]
    \color{white}\bfseries\large \thepartiecours{} --- #1%
    \signet[partie-\arabic{partiecours}]{1}{\thepartiecours{} --- #1}
  \end{tcolorbox}%
  \addcontentsline{toc}{section}{\thepartiecours{} --- #1}%
}

\newtcolorbox{boitecours}[3][]{%
  enhanced,breakable,
  colback=white,colframe=#2,coltitle=white,
  fonttitle=\bfseries,
  attach boxed title to top left={xshift=8pt,yshift=-\tcboxedtitleheight/2},
  boxed title style={colback=#2,arc=2pt,boxrule=0pt},
  boxrule=0.9pt,arc=2pt,
  left=8pt,right=8pt,top=8pt,bottom=6pt,
  title={#3},#1}

\newcounter{methode}
\newenvironment{definitionbox}[1][Définition]%
  {\begin{boitecours}[phantom={\signet{2}{#1}}]{coulDef}{#1}}{\end{boitecours}}
\newenvironment{proprietebox}[1][Propriété]%
  {\begin{boitecours}[phantom={\signet{2}{#1}}]{coulProp}{#1}}{\end{boitecours}}
\newenvironment{methodebox}[1][Méthode]{\stepcounter{methode}%
  \begin{boitecours}[phantom={\signet[methode-\arabic{methode}]{2}{#1}}]{coulMeth}{#1}}%
  {\end{boitecours}}

\newcommand{\etape}[1]{\par\smallskip\textbf{\color{coulMeth}#1}\par\nobreak\smallskip}
\newcommand{\enonce}[1]{\emph{#1}\par\smallskip}
\newcommand{\reponse}[1]{\par\smallskip
  {\footnotesize\itshape\color{black!55}Réponses : #1}}
\newcommand{\demo}[1]{\textbf{\color{coulMeth}Démonstration.} #1}
\newcommand{\afaire}[1]{%
  \par\smallskip
  \noindent{\small\color{coulExo}$\blacktriangleright$\ \textbf{À faire :}
    fiche d'exercices du chapitre, #1.}\par\smallskip}

\newenvironment{calculs}{\setlength{\jot}{5pt}\@fleqntrue\@mathmargin\retraitcalculs
  \setlength{\abovedisplayskip}{4pt}\setlength{\belowdisplayskip}{4pt}%
  \setlength{\abovedisplayshortskip}{4pt}\setlength{\belowdisplayshortskip}{4pt}%
  \csname alignat*\endcsname{2}}%
  {\csname endalignat*\endcsname}
\newcommand{\rem}[1]{\text{\small\itshape\color{black!60}#1}}
\newcommand{\explique}[2]{\par\smallskip\noindent
  \begin{minipage}[t]{0.57\linewidth}\raggedright #1\end{minipage}\hfill
  \begin{minipage}[t]{0.40\linewidth}\raggedright\leavevmode
    {\small\itshape\color{black!60}#2}\end{minipage}\par}
\newcommand{\cas}[1]{\par\smallskip\noindent\textbf{\color{coulExo}#1}\nobreak}
\newcommand{\calc}[1]{\par\smallskip\textbf{\color{coulExo}#1}\par\nobreak}

\newtcolorbox{remarquebox}[1][Remarque]{%
  enhanced,breakable,colback=white,colframe=coulRem,
  boxrule=0pt,leftrule=2.5pt,arc=0pt,outer arc=0pt,
  left=8pt,right=6pt,top=5pt,bottom=5pt,
  before upper={\textbf{\color{coulRem}#1}\par\smallskip}}

\newtcolorbox{exemplebox}[1][Exemple]{%
  enhanced,breakable,colback=white,colframe=coulEx,
  boxrule=0pt,leftrule=2.5pt,arc=0pt,outer arc=0pt,
  left=8pt,right=6pt,top=5pt,bottom=5pt,
  before upper={\textbf{\color{coulEx}#1}\par\smallskip}}

\pgfplotsset{
  repere/.style={
    axis lines=middle,
    axis line style={-{Stealth[length=2.4mm]},thick},
    every axis x label/.style={at={(ticklabel* cs:1.0)},anchor=west},
    every axis y label/.style={at={(ticklabel* cs:1.0)},anchor=south},
    label style={font=\footnotesize},
    tick label style={font=\scriptsize},
    grid=both,
    major grid style={grille},
    minor tick num=0,
    tick align=inside,
    clip mode=individual,
  },
}
\tikzset{
  courbe/.style={coulCourbe,line width=1.1pt,smooth},
  courbeB/.style={coulCourbeB,line width=1.1pt,smooth},
  courbeC/.style={coulCourbeC,line width=1.1pt,smooth},
}

\newlist{etapes}{enumerate}{1}
\setlist[etapes]{label=\textbf{\arabic*.},leftmargin=*,itemsep=1pt,topsep=2pt}

\newlist{expressions}{enumerate}{1}
\setlist[expressions]{label={\Alph*\ $=$},leftmargin=*,itemsep=3pt,topsep=3pt}

\newcounter{exo}
\renewcommand{\theexo}{\ifnum\value{exo}<10 0\fi\arabic{exo}}
\newtcolorbox{exercice}[1][]{%
  enhanced,breakable,
  colback=white,colframe=coulExo,
  boxrule=0.9pt,arc=2pt,
  left=8pt,right=8pt,top=8pt,bottom=6pt,
  coltitle=white,fonttitle=\bfseries,
  attach boxed title to top left={xshift=8pt,yshift=-\tcboxedtitleheight/2},
  boxed title style={colback=coulExo,arc=2pt,boxrule=0pt},
  phantom={\signet{2}{Exercice \arabic{exo}}},
  title={Exercice \theexo\ifx\relax#1\relax\else{} --- #1\fi}}
\newenvironment{exo}[1][\relax]{\refstepcounter{exo}\begin{exercice}[#1]}{\end{exercice}}

  \newenvironment{correction}{\par}{\par}
  \newcommand{\autableau}{}
  \newcommand{\etapeexemples}{\etape{Exemples corrigés}}
  \newcommand{\etapetableau}[1]{\etape{#1}}
  \newenvironment{exempletableau}[1][Exemple]%
    {\begin{exemplebox}[#1]}{\end{exemplebox}}

\RequirePackage{comment}
\excludecomment{correctionavj}

\newcommand{\coursEntete}{}
\newcommand{\coursNiveau}{}
\newcommand{\coursTitre}{}
\newcommand{\coursSurTitre}{}
\newcommand{\coursSousTitre}{}

\newcommand{\enteteCours}[1]{\renewcommand{\coursEntete}{#1}}
\newcommand{\chapitrecours}[2]{\enteteCours{Chapitre #1 --- #2}}
\newcommand{\niveaucours}[1]{\renewcommand{\coursNiveau}{#1}}
\newcommand{\titrecours}[3][]{%
  \renewcommand{\coursTitre}{#2}%
  \renewcommand{\coursSurTitre}{#1}%
  \renewcommand{\coursSousTitre}{#3}}

\pagestyle{fancy}
\fancyhf{}
\fancyhead[L]{\footnotesize\color{coulPart}\textbf{\coursEntete}}
\fancyhead[R]{\footnotesize\color{coulPart}\textbf{\coursNiveau}}
\fancyfoot[C]{\footnotesize\thepage}
\renewcommand{\headrulewidth}{0.4pt}
\renewcommand{\headrule}{\hbox to\headwidth{\color{coulPart!40}\leaders\hrule height \headrulewidth\hfill}}

\setlength{\parskip}{2pt}

\AtBeginDocument{%
  \begin{center}
    \begin{tcolorbox}[enhanced,width=0.72\linewidth,colback=white,colframe=coulPart,
        boxrule=1.6pt,arc=3pt,halign=center,top=10pt,bottom=10pt]
      \ifx\coursSurTitre\empty
        {\Huge\bfseries\color{coulPart} \coursTitre}\\[4pt]
      \else
        {\Huge\bfseries\color{coulPart} \coursTitre}\\[3pt]
        {\Large\bfseries\color{coulPart} \coursSurTitre}\\[5pt]
      \fi
      {\normalsize \coursSousTitre}%
      
    \end{tcolorbox}
  \end{center}%
}
\makeatother

\chapitrecours{5}{Suites numériques}
\niveaucours{1\textsuperscript{re} spécialité mathématiques}
\titrecours{SUITES NUMÉRIQUES}{Cours --- Première spécialité}

\begin{document}

\newtcblisting{python}{%
  listing only,enhanced,breakable,
  colback=black!3,colframe=black!35,boxrule=0.5pt,arc=1pt,
  left=6pt,right=6pt,top=2pt,bottom=2pt,
  listing options={language=Python,basicstyle=\ttfamily\small,
    keywordstyle=\color{coulMeth}\bfseries,columns=fullflexible,
    keepspaces=true,showstringspaces=false,aboveskip=0pt,belowskip=0pt}}
\newcommand{\reperesuite}[4]{%
  \draw[grille,step=1] (-1,#2) grid (#1,#3);
  \draw[-{Stealth[length=2mm]},thick] (-1,0)--({#1+0.4},0)
    node[below left=-1pt,font=\footnotesize]{$n$};
  \draw[-{Stealth[length=2mm]},thick] (0,#2)--(0,{#3+0.4})
    node[below left=-1pt,font=\footnotesize]{$#4$};
  \node[below left,font=\scriptsize,inner sep=1.5pt] at (0,0) {$0$};}

\begin{remarquebox}[Comment utiliser ce cours]
Chaque partie commence par ce qu'il faut \textbf{savoir} (définitions, propriétés), puis
vient ce qu'il faut \textbf{savoir faire} : une méthode, avec des
\textbf{exemples corrigés} faits en classe, puis des exercices
\og\textbf{À vous de jouer}\fg{} à chercher seul.

\end{remarquebox}

\partiecours{Définition d'une suite}

\begin{definitionbox}[Définition 1 --- Suite numérique]
\begin{minipage}[c]{0.56\linewidth}\raggedright
Une \textbf{suite} $(u_{n})$ associe à chaque entier naturel $n$ (à partir d'un
certain rang $n_{0}$) un nombre réel noté $u_{n}$.

\smallskip
C'est une fonction définie sur les entiers : on écrit $u_{n}$ plutôt que $u(n)$.
\begin{itemize}
  \item $n$ est le \textbf{rang} (ou l'indice) ;
  \item $u_{n}$ est le \textbf{terme de rang $n$} ;
  \item $(u_{n})$, avec les parenthèses, désigne toute la suite.
\end{itemize}
\end{minipage}\hfill
\begin{minipage}[c]{0.40\linewidth}\centering
\renewcommand{\arraystretch}{1.4}
\begin{tabular}{|c|c|c|c|c|c|}
\hline
rang $n$ & $0$ & $1$ & $2$ & $3$ & $4$\\\hline
terme $u_{n}$ & $0$ & $1$ & $4$ & $9$ & $16$\\\hline
\end{tabular}

\smallskip
{\footnotesize la suite des carrés : $u_{3}=9$}
\end{minipage}
\end{definitionbox}

\begin{remarquebox}[Termes consécutifs]
$u_{n}$ et $u_{n+1}$ se suivent : $u_{n+1}$ est le terme \textbf{suivant} $u_{n}$ et
$u_{n-1}$ est le terme \textbf{précédent}. Si le premier terme est $u_{0}$, le
5\textsuperscript{e} terme est $u_{4}$ : on compte à partir de $0$.
\end{remarquebox}

\begin{remarquebox}[L'indice s'écrit en petit]
Dans $u_{4}$, le rang $4$ est un \textbf{indice} : il s'écrit \textbf{en petit, en bas à
droite} du $u$. Écrit à la même taille, \og $u\,4$\fg{} se lirait $u\times 4$ : ce n'est
pas $u_{4}$. De même, $u_{n+1}$ (terme suivant) n'est pas $u_{n}+1$ (terme plus un).
\end{remarquebox}

\begin{methodebox}[Méthode 1 --- Repérer les termes et les rangs]
\etapeexemples
Une suite $(u_{n})$, de premier terme $u_{0}$, commence par :
\quad $3$ ;\; $8$ ;\; $15$ ;\; $24$ ;\; $35$ ;\; $48$.
\begin{questions}[label=\alph*)]
  \item Donner $u_{0}$ et $u_{4}$.
  \item Quel est le rang du terme $24$ ?
  \item Quel est le 3\textsuperscript{e} terme de la suite ?
  \item Donner le terme qui suit $u_{4}$, puis celui qui le précède.
\end{questions}
\begin{correction}
\explique{On numérote les termes à partir de $0$ :}{le premier terme est $u_{0}$}
\begin{center}
\renewcommand{\arraystretch}{1.3}
\begin{tabular}{|c|c|c|c|c|c|c|}
\hline
$n$ & $0$ & $1$ & $2$ & $3$ & $4$ & $5$\\\hline
$u_{n}$ & $3$ & $8$ & $15$ & $24$ & $35$ & $48$\\\hline
\end{tabular}
\end{center}
\cas{a)}
\explique{$u_{0}=3$ et $u_{4}=35$.}{on lit sous $n=0$ et sous $n=4$}
\cas{b)}
\explique{$24=u_{3}$ : le terme $24$ est de rang $3$.}{le rang est l'indice}
\cas{c)}
\explique{Le 3\textsuperscript{e} terme est $u_{2}=15$.}{on compte $u_{0}$, $u_{1}$
puis $u_{2}$}
\cas{d)}
\explique{Le terme suivant est $u_{5}=48$ ; le précédent est $u_{3}=24$.}{$u_{4+1}$
et $u_{4-1}$}
\end{correction}

\etape{À vous de jouer}
Une suite $(v_{n})$, de premier terme $v_{0}$, commence par :
\quad $2$ ;\; $6$ ;\; $18$ ;\; $54$ ;\; $162$ ;\; $486$.
\begin{questions}[label=\alph*)]
  \item Donner $v_{0}$ et $v_{4}$.
  \item Quel est le rang du terme $54$ ?
  \item Quel est le 3\textsuperscript{e} terme de la suite ?
  \item Donner le terme qui suit $v_{4}$, puis celui qui le précède.
\end{questions}
\reponse{a) $v_{0}=2$ et $v_{4}=162$.\; b) Rang $3$.\; c) $v_{2}=18$.\;
d) $v_{5}=486$ et $v_{3}=54$.}
\end{methodebox}

\begin{methodebox}[Méthode 2 --- Trouver le premier rang d'une suite]
\etapeexemples
Pour chaque suite, donner le premier rang $n_{0}$ à partir duquel elle est définie.
\begin{questions}[label=\alph*)]
  \item $u_{n}=\dfrac{2}{n}$
  \item $v_{n}=\sqrt{n-4}$
  \item $w_{n}=\dfrac{3}{\sqrt{n-1}}$
\end{questions}
\begin{correction}
\cas{a)}
\explique{On ne divise pas par $0$ : il faut $n\neq 0$.}{valeur interdite}
\explique{$(u_{n})$ est définie à partir du rang $n_{0}=1$.}{premier entier
différent de $0$}
\cas{b)}
\explique{Sous la racine, il faut $n-4\ge 0$, soit $n\ge 4$.}{une racine carrée
n'existe que pour un nombre positif}
\explique{$(v_{n})$ est définie à partir du rang $n_{0}=4$.}{$v_{4}=\sqrt{0}=0$}
\cas{c)}
\explique{Il faut $n-1\ge 0$ (racine) et $\sqrt{n-1}\neq 0$ (division), soit
$n-1>0$.}{les deux conditions}
\explique{$n>1$ : $(w_{n})$ est définie à partir du rang $n_{0}=2$.}{$n$ est
entier : après $1$ vient $2$}
\end{correction}

\etape{À vous de jouer}
Pour chaque suite, donner le premier rang $n_{0}$ à partir duquel elle est définie.
\begin{questions}[label=\alph*)]
  \item $a_{n}=-\dfrac{4}{n}$
  \item $b_{n}=\sqrt{n-6}$
  \item $c_{n}=\dfrac{5}{\sqrt{n-2}}$
\end{questions}
\reponse{a) $n_{0}=1$.\; b) $n_{0}=6$.\; c) $n_{0}=3$.}
\end{methodebox}

\afaire{exercices 1 à 4}

\partiecours{Suites définies explicitement}

\begin{definitionbox}[Définition 2 --- Suite définie explicitement]
Une suite est \textbf{définie explicitement} lorsque $u_{n}$ s'écrit directement en
fonction de $n$ : $u_{n}=f(n)$. L'expression de $u_{n}$ s'appelle le
\textbf{terme général}.

\smallskip
Exemple : avec $u_{n}=2n+1$, on obtient $u_{10}=21$ sans calculer les termes
précédents.
\end{definitionbox}

\begin{methodebox}[Méthode 3 --- Calculer un terme d'une suite explicite]
\etapeexemples
Pour tout entier naturel $n$, on pose $u_{n}=n^{2}-3n+5$ et
$v_{n}=\dfrac{2n+1}{n+2}$.
\begin{questions}[label=\alph*)]
  \item Calculer $u_{0}$ et $u_{4}$.
  \item Calculer $v_{0}$ et $v_{3}$.
\end{questions}
\begin{correction}
\cas{a)}
\begin{calculs}
u_{0} &= 0^{2}-3\times 0+5 &\qquad& \rem{on remplace $n$ par $0$}\\
u_{0} &= 5 &&\\[4pt]
u_{4} &= 4^{2}-3\times 4+5 &\qquad& \rem{on remplace $n$ par $4$}\\
u_{4} &= 16-12+5 = 9 &&
\end{calculs}
\cas{b)}
\begin{calculs}
v_{0} &= \frac{2\times 0+1}{0+2} &\qquad& \rem{on remplace $n$ par $0$}\\
v_{0} &= \frac{1}{2} &&\\[4pt]
v_{3} &= \frac{2\times 3+1}{3+2} &\qquad& \rem{on remplace $n$ par $3$}\\
v_{3} &= \frac{7}{5} = 1{,}4 &&
\end{calculs}
\end{correction}

\etape{À vous de jouer}
Pour tout entier naturel $n$, on pose $a_{n}=2n^{2}-n+1$ et
$b_{n}=\dfrac{3n-1}{n+1}$.
\begin{questions}[label=\alph*)]
  \item Calculer $a_{0}$ et $a_{3}$.
  \item Calculer $b_{0}$ et $b_{4}$.
\end{questions}
\reponse{a) $a_{0}=1$ et $a_{3}=16$.\; b) $b_{0}=-1$ et $b_{4}=\frac{11}{5}=2{,}2$.}
\end{methodebox}

\begin{remarquebox}[Attention : $u_{n+1}$ n'est pas $u_{n}+1$]
Pour obtenir $u_{n+1}$, on remplace $n$ par $(n+1)$ \textbf{dans l'expression} ;
$u_{n}+1$, c'est le terme $u_{n}$ auquel on ajoute $1$. De même, $u_{n-1}$ s'obtient en
remplaçant $n$ par $(n-1)$.
\end{remarquebox}

\begin{methodebox}[Méthode 4 --- Exprimer $u_{n+1}$ et $u_{n-1}$ en fonction de $n$]
\etapeexemples
Pour tout entier naturel $n$, on pose $u_{n}=n^{2}+3$ et $v_{n}=\dfrac{n}{n+2}$.
\begin{questions}[label=\alph*)]
  \item Exprimer $u_{n+1}$ en fonction de $n$, puis le comparer à $u_{n}+1$.
  \item Exprimer $u_{n-1}$ en fonction de $n$ (pour $n\ge 1$).
  \item Exprimer $v_{n+1}$ en fonction de $n$.
\end{questions}
\begin{correction}
\cas{a)}
\begin{calculs}
u_{n+1} &= (n+1)^{2}+3 &\qquad& \rem{on remplace $n$ par $(n+1)$}\\
u_{n+1} &= n^{2}+2n+1+3 &\qquad& \rem{$(a+b)^{2}=a^{2}+2ab+b^{2}$}\\
u_{n+1} &= n^{2}+2n+4 &&\\[4pt]
u_{n}+1 &= n^{2}+3+1 &\qquad& \rem{on ajoute $1$ au terme $u_{n}$}\\
u_{n}+1 &= n^{2}+4 &&
\end{calculs}
\explique{Les deux expressions sont différentes : $u_{n+1}\neq u_{n}+1$.}{elles
diffèrent de $2n$}
\cas{b)}
\begin{calculs}
u_{n-1} &= (n-1)^{2}+3 &\qquad& \rem{on remplace $n$ par $(n-1)$}\\
u_{n-1} &= n^{2}-2n+1+3 &\qquad& \rem{$(a-b)^{2}=a^{2}-2ab+b^{2}$}\\
u_{n-1} &= n^{2}-2n+4 &&
\end{calculs}
\cas{c)}
\begin{calculs}
v_{n+1} &= \frac{n+1}{(n+1)+2} &\qquad& \rem{on remplace chaque $n$ par $(n+1)$}\\
v_{n+1} &= \frac{n+1}{n+3} &&
\end{calculs}
\end{correction}

\etape{À vous de jouer}
Pour tout entier naturel $n$, on pose $a_{n}=n^{2}+n$ et $b_{n}=\dfrac{2n}{n+1}$.
\begin{questions}[label=\alph*)]
  \item Exprimer $a_{n+1}$ en fonction de $n$, puis le comparer à $a_{n}+1$.
  \item Exprimer $a_{n-1}$ en fonction de $n$ (pour $n\ge 1$).
  \item Exprimer $b_{n+1}$ en fonction de $n$.
\end{questions}
\reponse{a) $a_{n+1}=n^{2}+3n+2$ et $a_{n}+1=n^{2}+n+1$ : différents.\;
b) $a_{n-1}=n^{2}-n$.\; c) $b_{n+1}=\frac{2n+2}{n+2}$.}
\end{methodebox}

\afaire{exercices 5 à 8}

\partiecours{Suites définies par récurrence}

\begin{definitionbox}[Définition 3 --- Suite définie par récurrence]
\begin{minipage}[c]{0.56\linewidth}\raggedright
Une suite est \textbf{définie par récurrence} lorsqu'on donne :
\begin{itemize}
  \item son premier terme, par exemple $u_{0}=2$ ;
  \item une \textbf{relation de récurrence}, qui donne chaque terme à partir du
        précédent, par exemple $u_{n+1}=3u_{n}-1$.
\end{itemize}
On calcule alors les termes \textbf{de proche en proche} : pour connaître $u_{20}$, il
faut d'abord $u_{19}$, donc $u_{18}$\dots
\end{minipage}\hfill
\begin{minipage}[c]{0.40\linewidth}\centering
\begin{tikzpicture}[>={Stealth[length=2mm]}]
  \node (a) at (0,0) {$u_{0}=2$};
  \node (b) at (2.1,0) {$u_{1}=5$};
  \node (c) at (4.2,0) {$u_{2}=14$};
  \draw[->,coulMeth,thick] (a.north east) to[bend left=40]
    node[above,font=\scriptsize]{$\times3-1$} (b.north west);
  \draw[->,coulMeth,thick] (b.north east) to[bend left=40]
    node[above,font=\scriptsize]{$\times3-1$} (c.north west);
\end{tikzpicture}
\end{minipage}
\end{definitionbox}

\begin{methodebox}[Méthode 5 --- Calculer des termes de proche en proche]
\etapeexemples
Pour chaque suite, calculer les trois termes qui suivent le premier.
\begin{questions}[label=\alph*)]
  \item $u_{0}=4$ et, pour tout entier naturel $n$, $u_{n+1}=2u_{n}-3$.
  \item $v_{0}=1$ et, pour tout entier naturel $n$, $v_{n+1}=v_{n}+2n+1$.
  \item $w_{1}=10$ et, pour tout entier $n\ge 1$, $w_{n+1}=\dfrac{w_{n}}{2}+1$.
\end{questions}
\begin{correction}
\cas{a)}
\begin{calculs}
u_{1} &= 2u_{0}-3 &\qquad& \rem{la relation avec $n=0$}\\
u_{1} &= 2\times 4-3 = 5 &\qquad& \rem{on remplace $u_{0}$ par sa valeur}\\[4pt]
u_{2} &= 2u_{1}-3 &\qquad& \rem{la relation avec $n=1$}\\
u_{2} &= 2\times 5-3 = 7 &&\\[4pt]
u_{3} &= 2u_{2}-3 &&\\
u_{3} &= 2\times 7-3 = 11 &&
\end{calculs}
\cas{b)}
\begin{calculs}
v_{1} &= v_{0}+2\times 0+1 &\qquad& \rem{$n=0$ : on remplace aussi $n$ par $0$}\\
v_{1} &= 1+0+1 = 2 &&\\[4pt]
v_{2} &= v_{1}+2\times 1+1 &\qquad& \rem{$n=1$ : le rang du terme précédent}\\
v_{2} &= 2+2+1 = 5 &&\\[4pt]
v_{3} &= v_{2}+2\times 2+1 &\qquad& \rem{$n=2$}\\
v_{3} &= 5+4+1 = 10 &&
\end{calculs}
\cas{c)}
\begin{calculs}
w_{2} &= \frac{w_{1}}{2}+1 &\qquad& \rem{le premier terme est $w_{1}$ : on part de $n=1$}\\
w_{2} &= \frac{10}{2}+1 = 6 &&\\[4pt]
w_{3} &= \frac{6}{2}+1 = 4 &&\\
w_{4} &= \frac{4}{2}+1 = 3 &&
\end{calculs}
\end{correction}

\etape{À vous de jouer}
Pour chaque suite, calculer les trois termes qui suivent le premier.
\begin{questions}[label=\alph*)]
  \item $a_{0}=3$ et, pour tout entier naturel $n$, $a_{n+1}=3a_{n}-5$.
  \item $b_{0}=2$ et, pour tout entier naturel $n$, $b_{n+1}=b_{n}+4n-1$.
  \item $c_{1}=20$ et, pour tout entier $n\ge 1$, $c_{n+1}=\dfrac{c_{n}}{2}+2$.
\end{questions}
\reponse{a) $a_{1}=4$ ; $a_{2}=7$ ; $a_{3}=16$.\; b) $b_{1}=1$ ; $b_{2}=4$ ;
$b_{3}=11$.\; c) $c_{2}=12$ ; $c_{3}=8$ ; $c_{4}=6$.}
\end{methodebox}

\begin{methodebox}[Méthode 6 --- Traduire une situation par une relation de récurrence]
\etapeexemples
Une médiathèque compte $5\,000$ abonnés au début de 2026. Chaque année, $20\,\%$ des
abonnés ne renouvellent pas leur abonnement et $1\,200$ nouveaux lecteurs s'inscrivent.
On note $u_{n}$ le nombre d'abonnés au début de l'année $2026+n$.
\begin{questions}[label=\alph*)]
  \item Donner $u_{0}$, puis calculer $u_{1}$.
  \item Justifier que, pour tout entier naturel $n$, $u_{n+1}=0{,}8u_{n}+1\,200$.
  \item Calculer $u_{2}$.
\end{questions}
\begin{correction}
\cas{a)}
\explique{$u_{0}=5\,000$.}{début de 2026 : $n=0$}
\begin{calculs}
u_{1} &= 5\,000-0{,}2\times 5\,000+1\,200 &\qquad& \rem{on retire $20\,\%$, on ajoute les nouveaux}\\
u_{1} &= 5\,000-1\,000+1\,200 = 5\,200 &&
\end{calculs}
\cas{b)}
\explique{Retirer $20\,\%$, c'est multiplier par $1-0{,}2=0{,}8$ : il reste
$0{,}8u_{n}$ abonnés.}{coefficient multiplicateur}
\explique{On ajoute les $1\,200$ inscrits : $u_{n+1}=0{,}8u_{n}+1\,200$.}{d'une
année à la suivante}
\cas{c)}
\begin{calculs}
u_{2} &= 0{,}8\times 5\,200+1\,200 &\qquad& \rem{la relation avec $n=1$}\\
u_{2} &= 4\,160+1\,200 = 5\,360 &&
\end{calculs}
\end{correction}

\etape{À vous de jouer}
Une citerne de jardin contient $2\,000$~L d'eau le 1\textsuperscript{er} juin. Chaque
semaine, on utilise $10\,\%$ de l'eau pour arroser et la pluie apporte $150$~L. On note
$v_{n}$ le volume d'eau, en litres, au bout de $n$ semaines.
\begin{questions}[label=\alph*)]
  \item Donner $v_{0}$, puis calculer $v_{1}$.
  \item Justifier que, pour tout entier naturel $n$, $v_{n+1}=0{,}9v_{n}+150$.
  \item Calculer $v_{2}$.
\end{questions}
\reponse{a) $v_{0}=2\,000$ et $v_{1}=1\,950$.\; b) Retirer $10\,\%$, c'est multiplier par
$0{,}9$, puis on ajoute $150$.\; c) $v_{2}=1\,905$.}
\end{methodebox}

\afaire{exercices 9 à 12}

\partiecours{Calculatrice}

\begin{proprietebox}[Application \texttt{Suites} de la NumWorks]
\begin{center}
\renewcommand{\arraystretch}{1.3}
\begin{tabular}{|>{\raggedright\arraybackslash}p{4.2cm}|>{\raggedright\arraybackslash}p{12.4cm}|}
\hline
\textbf{Pour} & \textbf{On fait}\\\hline
entrer une suite explicite & Onglet \texttt{Suites}, \texttt{Ajouter une suite},
  \texttt{Explicite}, puis par exemple $u_{n}=n^{2}-8n+20$ ($n$ avec la touche
  \fbox{\small$x,n,t$}).\\\hline
entrer une suite récurrente & \texttt{Ajouter une suite}, \texttt{Récurrente d'ordre 1},
  puis la relation $u_{n+1}=\dots$ ($u_{n}$ se trouve dans la boîte à outils
  \fbox{\small toolbox}) et le premier terme $u_{0}$.\\\hline
afficher des termes & Onglet \texttt{Tableau}, \texttt{Régler l'intervalle} :
  \texttt{N début}, \texttt{N fin}.\\\hline
représenter la suite & Onglet \texttt{Graphique} : les points $(n\,;u_{n})$ ; menu
  \texttt{Axes} pour régler la fenêtre.\\\hline
\end{tabular}
\end{center}
{\footnotesize\itshape Les intitulés peuvent varier selon la version du logiciel.}
\end{proprietebox}

\begin{methodebox}[Méthode 7 --- Calculer des termes avec la calculatrice]
\etapeexemples
Pour tout entier naturel $n$, on pose $u_{n}=n^{2}-8n+20$. La suite $(v_{n})$ est
définie par $v_{0}=1$ et $v_{n+1}=2v_{n}+3$.
\begin{questions}[label=\alph*)]
  \item Afficher les termes $u_{0}$ à $u_{8}$ et les recopier.
  \item Représenter $(u_{n})$ à la calculatrice.
  \item Donner $v_{10}$ et $v_{20}$.
\end{questions}
\begin{correction}
\cas{a)}
\explique{\texttt{Explicite} : $u_{n}=n^{2}-8n+20$ ; \texttt{Tableau} : \texttt{N début}
$0$, \texttt{N fin} $8$.}{du rang $0$ au rang $8$}
\begin{center}
\renewcommand{\arraystretch}{1.3}
\begin{tabular}{|c|*{9}{c|}}
\hline
$n$ & $0$ & $1$ & $2$ & $3$ & $4$ & $5$ & $6$ & $7$ & $8$\\\hline
$u_{n}$ & $20$ & $13$ & $8$ & $5$ & $4$ & $5$ & $8$ & $13$ & $20$\\\hline
\end{tabular}
\end{center}
\cas{b)}
\explique{\texttt{Graphique}, \texttt{Axes} : $X_{\min}=0$, $X_{\max}=9$, $Y_{\min}=0$,
$Y_{\max}=22$.}{un peu au-delà des valeurs du tableau}
\cas{c)}
\explique{\texttt{Récurrente d'ordre 1} : $v_{n+1}=2v_{n}+3$ et $v_{0}=1$.}{la relation
et le premier terme}
\explique{\texttt{Tableau} jusqu'à $20$ : $v_{10}=4\,093$ et
$v_{20}=4\,194\,301$.}{la calculatrice enchaîne les termes à notre place}
\end{correction}

\etape{À vous de jouer}
Pour tout entier naturel $n$, on pose $a_{n}=n^{2}-6n+11$. La suite $(b_{n})$ est
définie par $b_{0}=2$ et $b_{n+1}=3b_{n}-2$.
\begin{questions}[label=\alph*)]
  \item Afficher les termes $a_{0}$ à $a_{6}$ et les recopier.
  \item Représenter $(a_{n})$ à la calculatrice.
  \item Donner $b_{10}$ et $b_{15}$.
\end{questions}
\reponse{a) $11$ ; $6$ ; $3$ ; $2$ ; $3$ ; $6$ ; $11$.\; b) Par exemple $X$ de $0$ à
$7$ et $Y$ de $0$ à $12$.\; c) $b_{10}=59\,050$ et $b_{15}=14\,348\,908$.}
\end{methodebox}

\afaire{exercices 13 à 16}

\partiecours{Programmer en Python}

\begin{remarquebox}[Python sur la NumWorks]
Application \texttt{Python}, \texttt{Ajouter un script} : on tape la fonction, puis
\texttt{Exécuter le script} ; dans la console, on appelle la fonction, par exemple
\texttt{terme\_v(10)}. Une boucle \texttt{for i in range(n)} répète $n$ fois les lignes
décalées en dessous d'elle.
\end{remarquebox}

\begin{methodebox}[Méthode 8 --- Programmer le calcul d'un terme en Python]
\etapeexemples
On reprend la suite $(v_{n})$ définie par $v_{0}=1$ et $v_{n+1}=2v_{n}+3$.
\begin{questions}[label=\alph*)]
  \item Écrire une fonction \texttt{terme\_v(n)} qui renvoie $v_{n}$.
  \item Que renvoie \texttt{terme\_v(10)} ?
\end{questions}
\begin{correction}
\cas{a)}\par\nopagebreak
\begin{minipage}[t]{0.42\linewidth}\vspace{0pt}
\begin{python}
def terme_v(n):
    v = 1
    for i in range(n):
        v = 2*v + 3
    return v
\end{python}
\end{minipage}\hfill
\begin{minipage}[t]{0.54\linewidth}\vspace{0pt}\raggedright
{\small\itshape\color{black!60}\texttt{v = 1} donne $v_{0}$ ; la boucle tourne $n$ fois,
de $v_{0}$ à $v_{n}$ ; dans \texttt{v = 2*v + 3}, le nouveau \texttt{v} est le terme
suivant, l'ancien est le terme précédent}
\end{minipage}
\cas{b)}
\explique{\texttt{terme\_v(10)} renvoie \texttt{4093} : $v_{10}=4\,093$.}{comme avec
la calculatrice}
\end{correction}

\etape{À vous de jouer}
On reprend la suite $(b_{n})$ définie par $b_{0}=2$ et $b_{n+1}=3b_{n}-2$.
\begin{questions}[label=\alph*)]
  \item Écrire une fonction \texttt{terme\_b(n)} qui renvoie $b_{n}$.
  \item Que renvoie \texttt{terme\_b(10)} ?
\end{questions}
\reponse{a) \texttt{b = 2}, puis \texttt{b = 3*b - 2} dans la boucle.\;
b) \texttt{59050}.}
\end{methodebox}

\afaire{exercices 17 à 20}

\partiecours{Représentation graphique}

\begin{definitionbox}[Définition 4 --- Représentation graphique d'une suite]
\begin{minipage}[c]{0.56\linewidth}\raggedright
Dans un repère, la \textbf{représentation graphique} de la suite $(u_{n})$ est le nuage
des points de coordonnées $(n\,;u_{n})$.

\smallskip
Ces points sont \textbf{isolés} : $n$ ne prend que des valeurs entières, on ne les relie
donc pas.

\smallskip
{\footnotesize Ci-contre : $u_{n}=0{,}5n^{2}-3n+5$, de $u_{0}$ à $u_{6}$.}
\end{minipage}\hfill
\begin{minipage}[c]{0.40\linewidth}\centering
\begin{tikzpicture}[x=0.5cm,y=0.5cm]
  \reperesuite{7}{0}{6}{u_{n}}
  \gradx{1,2,3,4,5,6} \grady{1,2,3,4,5}
  \foreach \n/\u in {0/5,1/2.5,2/1,3/0.5,4/1,5/2.5,6/5}
    \path[coulCourbe] (\n,\u) node[croix]{};
\end{tikzpicture}
\end{minipage}
\end{definitionbox}

\begin{methodebox}[Méthode 9 --- Représenter une suite et conjecturer ses variations]
\etapeexemples
La suite $(u_{n})$ est définie par $u_{0}=1$ et, pour tout entier naturel $n$,
$u_{n+1}=0{,}5u_{n}+3$.
\begin{questions}[label=\alph*)]
  \item Dresser le tableau des termes $u_{0}$ à $u_{5}$ (arrondis à $0{,}01$).
  \item Représenter ces termes dans un repère.
  \item Conjecturer le sens de variation de la suite.
\end{questions}
\begin{correction}
\cas{a)}
\explique{$u_{1}=0{,}5\times 1+3=3{,}5$, puis de même pour les suivants.}{de proche en
proche, ou avec la calculatrice}
\begin{center}
\renewcommand{\arraystretch}{1.3}
\begin{tabular}{|c|*{6}{c|}}
\hline
$n$ & $0$ & $1$ & $2$ & $3$ & $4$ & $5$\\\hline
$u_{n}$ & $1$ & $3{,}5$ & $4{,}75$ & $5{,}38$ & $5{,}69$ & $5{,}84$\\\hline
\end{tabular}
\end{center}
\cas{b) et c)}\par\nopagebreak
\begin{minipage}[c]{0.58\linewidth}\raggedright
On place un point par colonne du tableau : $(0\,;1)$, $(1\,;3{,}5)$,
$(2\,;4{,}75)$\dots

\smallskip
Les points montent de gauche à droite : la suite semble \textbf{croissante}.
{\small\itshape\color{black!60}C'est une conjecture, démontrée à la partie VII.}
\end{minipage}\hfill
\begin{minipage}[c]{0.38\linewidth}\centering
\begin{tikzpicture}[x=0.55cm,y=0.55cm]
  \reperesuite{6}{0}{7}{u_{n}}
  \gradx{1,2,3,4,5} \grady{1,2,3,4,5,6}
  \foreach \n/\u in {0/1,1/3.5,2/4.75,3/5.375,4/5.6875,5/5.84375}
    \path[coulCourbe] (\n,\u) node[croix]{};
\end{tikzpicture}
\end{minipage}
\end{correction}

\etape{À vous de jouer}
La suite $(v_{n})$ est définie par $v_{0}=8$ et, pour tout entier naturel $n$,
$v_{n+1}=0{,}5v_{n}+1$.
\begin{questions}[label=\alph*)]
  \item Dresser le tableau des termes $v_{0}$ à $v_{5}$ (arrondis à $0{,}01$).
  \item Représenter ces termes dans un repère.
  \item Conjecturer le sens de variation de la suite.
\end{questions}
\reponse{a) $8$ ; $5$ ; $3{,}5$ ; $2{,}75$ ; $2{,}38$ ; $2{,}19$.\; c) Les points
descendent : la suite semble décroissante.}
\end{methodebox}

\afaire{exercices 21 à 24}

\partiecours{Sens de variation : signe de $u_{n+1}-u_{n}$}

\begin{definitionbox}[Définition 5 --- Suite croissante, suite décroissante]
\begin{minipage}[c]{0.60\linewidth}\raggedright
La suite $(u_{n})$ est :
\begin{itemize}
  \item \textbf{croissante} si, pour tout $n$, $u_{n+1}\ge u_{n}$ ;
  \item \textbf{décroissante} si, pour tout $n$, $u_{n+1}\le u_{n}$ ;
  \item \textbf{constante} si, pour tout $n$, $u_{n+1}=u_{n}$.
\end{itemize}
Avec $>$ ou $<$, elle est \textbf{strictement} croissante ou décroissante. Si l'inégalité
n'est vraie que pour $n\ge p$, la suite est croissante (ou décroissante) \textbf{à
partir du rang $p$}.
\end{minipage}\hfill
\begin{minipage}[c]{0.38\linewidth}\centering
\begin{tikzpicture}[x=0.35cm,y=0.35cm]
  \reperesuite{5}{0}{5}{u_{n}}
  \foreach \n/\u in {0/1,1/1.5,2/2.5,3/3,4/4}
    \path[coulCourbe] (\n,\u) node[croix]{};
  \node[font=\scriptsize] at (2.5,-1.1) {croissante};
\end{tikzpicture}\hspace{4mm}%
\begin{tikzpicture}[x=0.35cm,y=0.35cm]
  \reperesuite{5}{0}{5}{u_{n}}
  \foreach \n/\u in {0/4,1/3,2/2.5,3/1.5,4/1}
    \path[coulCourbe] (\n,\u) node[croix]{};
  \node[font=\scriptsize] at (2.5,-1.1) {décroissante};
\end{tikzpicture}
\end{minipage}
\end{definitionbox}

\begin{proprietebox}[Propriété 1 --- Signe de la différence]
\begin{itemize}
  \item Si, pour tout $n$, $u_{n+1}-u_{n}\ge 0$, alors $(u_{n})$ est croissante.
  \item Si, pour tout $n$, $u_{n+1}-u_{n}\le 0$, alors $(u_{n})$ est décroissante.
\end{itemize}
En effet, $u_{n+1}-u_{n}\ge 0$ revient à $u_{n+1}\ge u_{n}$ (on ajoute $u_{n}$ aux deux
membres).
\end{proprietebox}

\begin{methodebox}[Méthode 10 --- Étudier le sens de variation par la différence]
\etapeexemples
Étudier le sens de variation de chaque suite.
\begin{questions}[label=\alph*)]
  \item $u_{n}=2n^{2}+n$ pour tout entier naturel $n$.
  \item $v_{0}=3$ et, pour tout entier naturel $n$, $v_{n+1}=v_{n}-n^{2}-1$.
  \item $w_{n}=n^{2}-6n$ pour tout entier naturel $n$.
\end{questions}
\begin{correction}
\cas{a)}
\explique{On exprime d'abord $u_{n+1}$ :}{on remplace $n$ par $(n+1)$}
\begin{calculs}
u_{n+1} &= 2(n+1)^{2}+(n+1) &\qquad& \rem{on recopie l'expression}\\
u_{n+1} &= 2n^{2}+4n+2+n+1 &\qquad& \rem{$(n+1)^{2}=n^{2}+2n+1$}\\
u_{n+1} &= 2n^{2}+5n+3 &&
\end{calculs}
\explique{Puis la différence :}{on soustrait $u_{n}$ entre parenthèses}
\begin{calculs}
u_{n+1}-u_{n} &= 2n^{2}+5n+3-(2n^{2}+n) &&\\
u_{n+1}-u_{n} &= 2n^{2}+5n+3-2n^{2}-n &\qquad& \rem{le signe $-$ change tous les signes}\\
u_{n+1}-u_{n} &= 4n+3 &&
\end{calculs}
\explique{Comme $n\ge 0$, on a $4n+3>0$ : la suite $(u_{n})$ est strictement
croissante.}{une différence positive pour tout $n$}
\cas{b)}
\begin{calculs}
v_{n+1}-v_{n} &= v_{n}-n^{2}-1-v_{n} &\qquad& \rem{on remplace $v_{n+1}$ par la relation}\\
v_{n+1}-v_{n} &= -n^{2}-1 &&
\end{calculs}
\explique{$n^{2}\ge 0$, donc $-n^{2}-1<0$ : la suite $(v_{n})$ est strictement
décroissante.}{pas besoin de $v_{0}$}
\cas{c)}
\begin{calculs}
w_{n+1} &= (n+1)^{2}-6(n+1) &\qquad& \rem{on recopie l'expression}\\
w_{n+1} &= n^{2}+2n+1-6n-6 &&\\
w_{n+1} &= n^{2}-4n-5 &&\\[4pt]
w_{n+1}-w_{n} &= n^{2}-4n-5-(n^{2}-6n) &&\\
w_{n+1}-w_{n} &= 2n-5 &\qquad& \rem{les $n^{2}$ s'éliminent}
\end{calculs}
\explique{Signe de $2n-5$ :}{une inéquation du premier degré}
\begin{calculs}
2n-5 &\ge 0 &&\\
2n &\ge 5 &&\\
n &\ge 2{,}5 &&
\end{calculs}
\explique{$n$ est entier : pour $n\ge 3$, $w_{n+1}-w_{n}>0$. La suite $(w_{n})$ est
strictement croissante à partir du rang $3$.}{avant, elle décroît : $w_{0}=0$ et
$w_{3}=-9$}
\end{correction}

\etape{À vous de jouer}
Étudier le sens de variation de chaque suite.
\begin{questions}[label=\alph*)]
  \item $a_{n}=3n^{2}+2n$ pour tout entier naturel $n$.
  \item $b_{0}=10$ et, pour tout entier naturel $n$, $b_{n+1}=b_{n}-2n^{2}-3$.
  \item $c_{n}=n^{2}-8n$ pour tout entier naturel $n$.
\end{questions}
\reponse{a) $a_{n+1}-a_{n}=6n+5>0$ : strictement croissante.\;
b) $b_{n+1}-b_{n}=-2n^{2}-3<0$ : strictement décroissante.\;
c) $c_{n+1}-c_{n}=2n-7$ : strictement croissante à partir du rang $4$.}
\end{methodebox}

\afaire{exercices 25 à 28}

\partiecours{Sens de variation : quotient $\dfrac{u_{n+1}}{u_{n}}$}

\noindent\textbf{D'où vient le quotient ?} Une suite est croissante quand
$u_{n+1}\ge u_{n}$ ; si ses termes sont strictement positifs, on peut diviser par
$u_{n}$.
\begin{calculs}
u_{n+1} &\ge u_{n} &\qquad& \rem{la suite est croissante (définition 5)}\\
\frac{u_{n+1}}{u_{n}} &\ge \frac{u_{n}}{u_{n}} &\qquad& \rem{on divise par
$u_{n}>0$ : le sens ne change pas}\\
\frac{u_{n+1}}{u_{n}} &\ge 1 &\qquad& \rem{$\dfrac{u_{n}}{u_{n}}=1$}
\end{calculs}
Comparer $u_{n+1}$ à $u_{n}$ revient donc à comparer le quotient
$\dfrac{u_{n+1}}{u_{n}}$ à $1$ ; de même, $u_{n+1}\le u_{n}$ donne
$\dfrac{u_{n+1}}{u_{n}}\le 1$.

\begin{proprietebox}[Propriété 2 --- Comparer le quotient à 1]
Si \textbf{tous les termes} de la suite $(u_{n})$ sont \textbf{strictement positifs} :
\begin{itemize}
  \item si, pour tout $n$, $\dfrac{u_{n+1}}{u_{n}}\ge 1$, alors $(u_{n})$ est
        croissante ;
  \item si, pour tout $n$, $\dfrac{u_{n+1}}{u_{n}}\le 1$, alors $(u_{n})$ est
        décroissante.
\end{itemize}
Cette méthode convient quand $u_{n}$ est un produit, en particulier avec une puissance
$3^{n}$, $0{,}5^{n}$\dots : le quotient se simplifie.
\end{proprietebox}

\begin{methodebox}[Méthode 11 --- Étudier le sens de variation par le quotient]
\etapeexemples
Étudier le sens de variation de chaque suite.
\begin{questions}[label=\alph*)]
  \item $u_{n}=5\times 3^{n}$ pour tout entier naturel $n$.
  \item $v_{n}=8\times 0{,}5^{n}$ pour tout entier naturel $n$.
  \item $w_{n}=\dfrac{2^{n}}{n}$ pour tout entier $n\ge 1$.
\end{questions}
\begin{correction}
\cas{a)}
\explique{$5>0$ et $3^{n}>0$, donc $u_{n}>0$ pour tout $n$.}{on vérifie d'abord le
signe des termes}
\begin{calculs}
\frac{u_{n+1}}{u_{n}} &= \frac{5\times 3^{n+1}}{5\times 3^{n}} &\qquad& \rem{on recopie}\\
\frac{u_{n+1}}{u_{n}} &= \frac{\barre{5}\times 3\times \barre{3^{n}}}{\barre{5}\times \barre{3^{n}}} &\qquad& \rem{$3^{n+1}=3\times 3^{n}$}\\
\frac{u_{n+1}}{u_{n}} &= 3 &&
\end{calculs}
\explique{$3>1$ : la suite $(u_{n})$ est strictement croissante.}{chaque terme est le
triple du précédent}
\cas{b)}
\explique{$8>0$ et $0{,}5^{n}>0$, donc $v_{n}>0$ pour tout $n$.}{signe des termes}
\begin{calculs}
\frac{v_{n+1}}{v_{n}} &= \frac{8\times 0{,}5^{n+1}}{8\times 0{,}5^{n}} &\qquad& \rem{on recopie}\\
\frac{v_{n+1}}{v_{n}} &= 0{,}5 &\qquad& \rem{même simplification qu'en a)}
\end{calculs}
\explique{$0{,}5<1$ : la suite $(v_{n})$ est strictement décroissante.}{chaque terme
est la moitié du précédent}
\cas{c)}
\explique{Pour $n\ge 1$, $2^{n}>0$ et $n>0$, donc $w_{n}>0$.}{signe des termes}
\begin{calculs}
\frac{w_{n+1}}{w_{n}} &= \frac{2^{n+1}}{n+1}\times\frac{n}{2^{n}} &\qquad& \rem{diviser par $w_{n}$, c'est multiplier par son inverse}\\
\frac{w_{n+1}}{w_{n}} &= \frac{2\times\barre{2^{n}}\times n}{(n+1)\times\barre{2^{n}}} &\qquad& \rem{$2^{n+1}=2\times 2^{n}$}\\
\frac{w_{n+1}}{w_{n}} &= \frac{2n}{n+1} &&
\end{calculs}
\explique{On compare ce quotient à $1$ par une différence :}{signe de
$\frac{w_{n+1}}{w_{n}}-1$}
\begin{calculs}
\frac{w_{n+1}}{w_{n}}-1 &= \frac{2n}{n+1}-\frac{n+1}{n+1} &\qquad& \rem{$1=\frac{n+1}{n+1}$}\\
\frac{w_{n+1}}{w_{n}}-1 &= \frac{2n-n-1}{n+1} &&\\
\frac{w_{n+1}}{w_{n}}-1 &= \frac{n-1}{n+1} &&
\end{calculs}
\explique{Pour $n\ge 1$ : $n-1\ge 0$ et $n+1>0$, donc $\frac{w_{n+1}}{w_{n}}\ge 1$. La
suite $(w_{n})$ est croissante.}{$w_{1}=w_{2}=2$, puis elle augmente}
\end{correction}

\etape{À vous de jouer}
Étudier le sens de variation de chaque suite.
\begin{questions}[label=\alph*)]
  \item $a_{n}=4\times 2^{n}$ pour tout entier naturel $n$.
  \item $b_{n}=6\times 0{,}8^{n}$ pour tout entier naturel $n$.
  \item $c_{n}=\dfrac{3^{n}}{n}$ pour tout entier $n\ge 1$.
\end{questions}
\reponse{a) Quotient $2>1$ : strictement croissante.\; b) Quotient $0{,}8<1$ :
strictement décroissante.\; c) Quotient $\frac{3n}{n+1}$ et
$\frac{3n}{n+1}-1=\frac{2n-1}{n+1}>0$ : strictement croissante.}
\end{methodebox}

\afaire{exercices 29 à 32}

\partiecours{Sens de variation : fonction associée}

\begin{proprietebox}[Propriété 3 --- Suite $u_{n}=f(n)$]
\begin{minipage}[c]{0.58\linewidth}\raggedright
Soit $(u_{n})$ définie par $u_{n}=f(n)$, où $f$ est une fonction définie sur
$\intfo{0}{+\infty}$.
\begin{itemize}
  \item Si $f$ est croissante sur $\intfo{0}{+\infty}$, alors $(u_{n})$ est croissante.
  \item Si $f$ est décroissante sur $\intfo{0}{+\infty}$, alors $(u_{n})$ est
        décroissante.
\end{itemize}
Les points $(n\,;u_{n})$ sont sur la courbe de $f$ : ils suivent ses variations.
\end{minipage}\hfill
\begin{minipage}[c]{0.38\linewidth}\centering
\begin{tikzpicture}[x=0.55cm,y=0.55cm]
  \reperesuite{7}{0}{4}{y}
  \gradx{1,2,3,4,5,6} \grady{1,2,3}
  \draw[coulCourbe,line width=1.1pt,domain=0:6.8,samples=60]
    plot (\x,{4-4/(\x+1)});
  \foreach \n in {0,1,...,6}
    \path[coulCourbeB] (\n,{4-4/(\n+1)}) node[croix]{};
  \node[coulCourbe,font=\footnotesize] at (5.6,2.9) {$\Cf$};
\end{tikzpicture}
\end{minipage}
\end{proprietebox}

\noindent\textbf{Démonstration} (cas croissant).
\begin{calculs}
n &< n+1 &\qquad& \rem{deux rangs qui se suivent}\\
f(n) &\le f(n+1) &\qquad& \rem{$f$ croissante conserve l'ordre}\\
u_{n} &\le u_{n+1} &\qquad& \rem{car $u_{n}=f(n)$ et $u_{n+1}=f(n+1)$}
\end{calculs}

\begin{remarquebox}[Attention : seulement pour une suite explicite]
Pour une suite définie par récurrence, $u_{n+1}=f(u_{n})$, les variations de $f$ ne
donnent pas celles de la suite : $x\mapsto 0{,}5x+1$ est croissante, pourtant la suite
$v_{0}=8$, $v_{n+1}=0{,}5v_{n}+1$ (méthode 9) est décroissante.
\end{remarquebox}

\begin{methodebox}[Méthode 12 --- Étudier le sens de variation à l'aide d'une fonction]
\etapeexemples
Étudier le sens de variation de chaque suite, définie pour tout entier naturel $n$.
\begin{questions}[label=\alph*)]
  \item $u_{n}=n^{3}+3n^{2}+2$
  \item $v_{n}=\dfrac{3n+1}{n+2}$
  \item $w_{n}=n^{3}-12n$
\end{questions}
\begin{correction}
\cas{a)}
\explique{On pose $f(x)=x^{3}+3x^{2}+2$ sur $\intfo{0}{+\infty}$ : $u_{n}=f(n)$.}{la
fonction associée}
\begin{calculs}
f'(x) &= 3x^{2}+6x &\qquad& \rem{on dérive}\\
f'(x) &= 3x(x+2) &\qquad& \rem{on factorise pour le signe}
\end{calculs}
\explique{Pour $x\ge 0$ : $3x\ge 0$ et $x+2>0$, donc $f'(x)\ge 0$ ; $f$ est croissante
sur $\intfo{0}{+\infty}$.}{$f'$ ne s'annule qu'en $0$}
\explique{La suite $(u_{n})$ est donc strictement croissante.}{propriété 3}
\cas{b)}
\explique{On pose $f(x)=\dfrac{3x+1}{x+2}$ sur $\intfo{0}{+\infty}$.}{la fonction
associée}
\begin{calculs}
f'(x) &= \frac{3(x+2)-(3x+1)\times 1}{(x+2)^{2}} &\qquad& \rem{$\left(\frac{u}{v}\right)'=\frac{u'v-uv'}{v^{2}}$}\\
f'(x) &= \frac{3x+6-3x-1}{(x+2)^{2}} &&\\
f'(x) &= \frac{5}{(x+2)^{2}} &&
\end{calculs}
\explique{$5>0$ et $(x+2)^{2}>0$ : $f'(x)>0$, donc $f$ est strictement croissante.}{un
carré non nul est positif}
\explique{La suite $(v_{n})$ est donc strictement croissante.}{propriété 3}
\cas{c)}
\explique{On pose $f(x)=x^{3}-12x$ sur $\intfo{0}{+\infty}$.}{la fonction associée}
\begin{calculs}
f'(x) &= 3x^{2}-12 &\qquad& \rem{on dérive}\\
f'(x) &= 3(x-2)(x+2) &\qquad& \rem{$x^{2}-4=(x-2)(x+2)$}\\
f(2) &= 8-24 = -16 &\qquad& \rem{image de la racine $2$}
\end{calculs}
\begin{center}
\begin{tikzpicture}[scale=0.9]
  \tkzTabInit[lgt=2,espcl=2.4]{$x$/0.8,$f'(x)$/0.8,$f(x)$/1.6}{$0$,$2$,$+\infty$}
  \tkzTabLine{,-,z,+,}
  \tkzTabVar{+/$0$,-/$-16$,+/{}}
\end{tikzpicture}
\end{center}
\explique{$f$ est croissante sur $\intfo{2}{+\infty}$ : la suite $(w_{n})$ est
croissante à partir du rang $2$.}{avant le rang $2$, elle décroît}
\end{correction}

\etape{À vous de jouer}
Étudier le sens de variation de chaque suite, définie pour tout entier naturel $n$.
\begin{questions}[label=\alph*)]
  \item $a_{n}=2n^{3}+6n^{2}-1$
  \item $b_{n}=\dfrac{2n+5}{n+1}$
  \item $c_{n}=n^{3}-27n$
\end{questions}
\reponse{a) $f'(x)=6x(x+2)\ge 0$ : strictement croissante.\;
b) $f'(x)=\frac{-3}{(x+1)^{2}}<0$ : strictement décroissante.\;
c) $f'(x)=3(x-3)(x+3)$ : croissante à partir du rang $3$.}
\end{methodebox}

\begin{remarquebox}[Quelle méthode choisir ?]
\begin{itemize}
  \item La \textbf{différence} $u_{n+1}-u_{n}$ : la méthode générale, la seule pour une
        suite définie par récurrence.
  \item Le \textbf{quotient} $\dfrac{u_{n+1}}{u_{n}}$ : termes strictement positifs,
        $u_{n}$ écrit comme un produit avec une puissance $n$.
  \item La \textbf{fonction associée} : suite explicite $u_{n}=f(n)$, quand $f$ se
        dérive facilement.
\end{itemize}
\end{remarquebox}

\afaire{exercices 33 à 36}

\partiecours{Notion de limite}

\begin{definitionbox}[Définition 6 --- Limite d'une suite]
\begin{minipage}[c]{0.58\linewidth}\raggedright
On s'intéresse aux termes $u_{n}$ quand $n$ devient de plus en plus grand.
\begin{itemize}
  \item La suite \textbf{tend vers $\ell$} (on dit aussi qu'elle \textbf{converge} vers
        $\ell$) si $u_{n}$ devient aussi proche de $\ell$ que l'on veut, pour $n$ assez
        grand. On note $\displaystyle\lim_{n\to+\infty}u_{n}=\ell$.
  \item La suite \textbf{tend vers $+\infty$} si $u_{n}$ devient aussi grand que l'on
        veut : $\displaystyle\lim_{n\to+\infty}u_{n}=+\infty$ (de même pour $-\infty$).
  \item Une suite peut n'avoir \textbf{aucune limite} : $(-1)^{n}$ vaut $1$ puis $-1$
        puis $1$\dots
\end{itemize}
\end{minipage}\hfill
\begin{minipage}[c]{0.38\linewidth}\centering
\begin{tikzpicture}[x=0.5cm,y=0.5cm]
  \reperesuite{9}{0}{4}{u_{n}}
  \gradx{1,2,3,4,5,6,7,8} \grady{1,2,3}
  \draw[coulCourbeB,dashed,thick] (0,3)--(9,3)
    node[above left,font=\footnotesize]{$\ell=3$};
  \foreach \n in {0,1,...,8}
    \path[coulCourbe] (\n,{3-2/(\n+1)}) node[croix]{};
\end{tikzpicture}

\smallskip
{\footnotesize $u_{n}=3-\dfrac{2}{n+1}$ : les points se rapprochent de la droite
d'équation $y=3$}
\end{minipage}

\medskip
Exemples à connaître :\quad $\displaystyle\lim_{n\to+\infty}\frac{1}{n}=0$ \;;\quad
$\displaystyle\lim_{n\to+\infty}n^{2}=+\infty$ \;;\quad
$\displaystyle\lim_{n\to+\infty}(-n)=-\infty$.
\end{definitionbox}

\begin{methodebox}[Méthode 13 --- Conjecturer la limite d'une suite]
\etapeexemples
Pour chaque suite, définie pour tout entier naturel $n$, calculer les termes indiqués,
puis conjecturer sa limite.
\begin{questions}[label=\alph*)]
  \item $u_{n}=3-\dfrac{2}{n+1}$ : $u_{9}$, $u_{99}$ et $u_{999}$.
  \item $v_{n}=n^{2}-5$ : $v_{10}$, $v_{100}$ et $v_{1\,000}$.
  \item $w_{n}=2\times(-1)^{n}$ : $w_{0}$ à $w_{3}$.
\end{questions}
\begin{correction}
\cas{a)}
\begin{calculs}
u_{9} &= 3-\frac{2}{10} = 2{,}8 &\qquad& \rem{on remplace $n$ par $9$}\\
u_{99} &= 3-\frac{2}{100} = 2{,}98 &&\\
u_{999} &= 3-\frac{2}{1\,000} = 2{,}998 &&
\end{calculs}
\explique{Les termes se rapprochent de $3$ : on conjecture
$\displaystyle\lim_{n\to+\infty}u_{n}=3$.}{$\frac{2}{n+1}$ devient très petit}
\cas{b)}
\begin{calculs}
v_{10} &= 10^{2}-5 = 95 &\qquad& \rem{on remplace $n$ par $10$}\\
v_{100} &= 100^{2}-5 = 9\,995 &&\\
v_{1\,000} &= 1\,000^{2}-5 = 999\,995 &&
\end{calculs}
\explique{Les termes deviennent aussi grands que l'on veut : on conjecture
$\displaystyle\lim_{n\to+\infty}v_{n}=+\infty$.}{$n^{2}$ grandit sans fin}
\cas{c)}
\explique{$w_{0}=2$ ; $w_{1}=-2$ ; $w_{2}=2$ ; $w_{3}=-2$.}{$(-1)^{n}$ vaut $1$ si $n$
est pair, $-1$ sinon}
\explique{Les termes alternent entre $2$ et $-2$ : la suite $(w_{n})$ n'a pas de
limite.}{ils ne se rapprochent d'aucun nombre}
\end{correction}

\etape{À vous de jouer}
Pour chaque suite, définie pour tout entier naturel $n$, calculer les termes indiqués,
puis conjecturer sa limite.
\begin{questions}[label=\alph*)]
  \item $a_{n}=5+\dfrac{4}{n+1}$ : $a_{9}$, $a_{99}$ et $a_{999}$.
  \item $b_{n}=2n^{2}+1$ : $b_{10}$, $b_{100}$ et $b_{1\,000}$.
  \item $c_{n}=3\times(-1)^{n}$ : $c_{0}$ à $c_{3}$.
\end{questions}
\reponse{a) $5{,}4$ ; $5{,}04$ ; $5{,}004$ : limite $5$.\; b) $201$ ; $20\,001$ ;
$2\,000\,001$ : limite $+\infty$.\; c) $3$ ; $-3$ ; $3$ ; $-3$ : pas de limite.}
\end{methodebox}

\afaire{exercices 37 à 40}

\partiecours{Problèmes : déterminer un seuil}

\begin{definitionbox}[Définition 7 --- Seuil]
Chercher un \textbf{seuil}, c'est chercher le \textbf{premier rang} $n$ à partir duquel
$u_{n}$ dépasse une valeur donnée (ou passe en dessous). On le trouve avec le tableau de
la calculatrice, ou avec un programme qui calcule les termes un par un \textbf{tant que}
la valeur n'est pas atteinte.
\end{definitionbox}

\begin{methodebox}[Méthode 14 --- Déterminer un seuil dans un problème]
\etapeexemples
Un club d'escalade compte $2\,000$ adhérents en septembre 2026. Chaque année, $10\,\%$
des adhérents ne se réinscrivent pas et $500$ nouveaux adhérents arrivent. On note
$u_{n}$ le nombre d'adhérents en septembre de l'année $2026+n$.
\begin{questions}[label=\alph*)]
  \item Calculer $u_{1}$ et $u_{2}$.
  \item Justifier que, pour tout entier naturel $n$, $u_{n+1}=0{,}9u_{n}+500$.
  \item Compléter la fonction ci-dessous pour qu'elle renvoie le nombre d'années au bout
        duquel le club dépasse $4\,000$ adhérents.
  \item À l'aide de la calculatrice, déterminer l'année où le club dépasse $4\,000$
        adhérents.
\end{questions}
\begin{python}
def seuil():
    u = 2000
    n = 0
    while .......... :
        u = ..........
        n = n + 1
    return n
\end{python}
\begin{correction}
\cas{a)}
\begin{calculs}
u_{1} &= 2\,000-0{,}1\times 2\,000+500 &\qquad& \rem{on retire $10\,\%$, on ajoute les nouveaux}\\
u_{1} &= 2\,000-200+500 = 2\,300 &&\\[4pt]
u_{2} &= 2\,300-230+500 = 2\,570 &\qquad& \rem{$10\,\%$ de $2\,300$ : $230$}
\end{calculs}
\cas{b)}
\explique{Retirer $10\,\%$, c'est multiplier par $0{,}9$ ; on ajoute ensuite les $500$
nouveaux : $u_{n+1}=0{,}9u_{n}+500$.}{comme à la méthode 6}
\cas{c)}\par\nopagebreak
\begin{minipage}[t]{0.42\linewidth}\vspace{0pt}
\begin{python}
def seuil():
    u = 2000
    n = 0
    while u <= 4000 :
        u = 0.9*u + 500
        n = n + 1
    return n
\end{python}
\end{minipage}\hfill
\begin{minipage}[t]{0.54\linewidth}\vspace{0pt}\raggedright
{\small\itshape\color{black!60}tant que le seuil n'est pas dépassé ($u\le 4\,000$, le
contraire de « dépasser »), on passe au terme suivant (relation de récurrence) et au rang
suivant ; la boucle s'arrête au premier terme qui dépasse $4\,000$, et \texttt{n} est
son rang}
\end{minipage}
\cas{d)}
\explique{\texttt{Tableau} de la calculatrice : $u_{10}\approx 3\,954$ et
$u_{11}\approx 4\,059$.}{le premier terme au-dessus de $4\,000$ est $u_{11}$}
\explique{\texttt{seuil()} renvoie $11$ : le club dépasse $4\,000$ adhérents en
septembre 2037.}{$2026+11=2037$}
\end{correction}

\etape{À vous de jouer}
Un étang contient $800$~kg d'algues le 1\textsuperscript{er} mai. Chaque semaine, on en
retire $20\,\%$ et elles repoussent de $60$~kg. On note $a_{n}$ la masse d'algues, en
kg, au bout de $n$ semaines.
\begin{questions}[label=\alph*)]
  \item Calculer $a_{1}$ et $a_{2}$.
  \item Justifier que, pour tout entier naturel $n$, $a_{n+1}=0{,}8a_{n}+60$.
  \item Compléter la fonction ci-dessous pour qu'elle renvoie le nombre de semaines au
        bout duquel la masse passe sous $400$~kg.
  \item À l'aide de la calculatrice, déterminer au bout de combien de semaines la masse
        passe sous $400$~kg.
\end{questions}
\begin{python}
def seuil():
    a = 800
    n = 0
    while .......... :
        a = ..........
        n = n + 1
    return n
\end{python}
\reponse{a) $a_{1}=700$ et $a_{2}=620$.\; b) Retirer $20\,\%$, c'est multiplier par
$0{,}8$, puis on ajoute $60$.\; c) \texttt{while a >= 400} et \texttt{a = 0.8*a + 60}.\;
d) $a_{7}\approx 404{,}9$ et $a_{8}\approx 383{,}9$ : au bout de $8$ semaines.}
\end{methodebox}
\afaire{exercices 41 à 47}

\end{document}
